\chapter{Переносная калибровка полного тензорного вихря}
\label{ch:v6-bosonic-defect-full-tensor-translation-calibration-gate}

\section{Проблема знака}

Предыдущий полный скрученный оператор не воспроизводил точную переносную
нулевую моду и одновременно давал две отрицательные пары. Это требовало
проверить не только собственные значения, но и действие гессиана на заранее
известную ковариантную касательную переноса. Такое отделение симметрийной
нулевой моды от внутренних возбуждений является обязательным в спектральной
теории вихрей \cite{translation-calibration-gustafson,translation-calibration-alonso}.

Переменная связности в полярном операторе была определена с противоположным
знаком относительно физической вариации $\delta A$. Поэтому фоновое
калибровочное условие в этих переменных имеет вид
\begin{equation}
 \mathcal G
 =\operatorname{div}\delta B
 +\frac{Z}{G}\langle R\Phi_0,\delta\Phi\rangle,
\end{equation}
а не с минусом перед материальным членом. Все остальные части оператора,
стационарный фон и граничные условия оставлены неизменными.

\section{Аналитическая переносная касательная}

В нейтральном характере $c=0$ и блоке $j=1$ построена касательная
\begin{equation}
 \delta\Phi_{+}
 =D_r\Phi_0+\frac{i}{r}D_\theta\Phi_0,
 \qquad
 \delta B_r=\frac{iK'}{r},
 \qquad
 \delta B_\theta=-\frac{K'}{r}.
\end{equation}
Она удовлетворяет сердцевинному условию
$\delta B_\theta=i\delta B_r$ и принадлежит той же дискретной области
определения, что и собственные векторы. Остаток её проекции на область не
превышает $3.38\cdot10^{-16}$.

На сетках $70,100,140,200,280$ нижний уровень равен
\begin{equation}
 0.00184869,\quad 0.000965984,\quad 0.000510129,\quad
 0.000254739,\quad 0.000131202.
\end{equation}
Рэлеевы отношения аналитической касательной равны соответственно
\begin{equation}
 0.00316195,\quad 0.00128654,\quad 0.000594204,\quad
 0.000274999,\quad 0.000136483.
\end{equation}
Перекрытие касательной с нижней собственной модой возрастает от
$0.9999619$ до $0.99999986$. Экстраполяция по $1/(N-1)^2$ даёт
\begin{equation}
 \lambda_{\mathrm{tr},\infty}=7.49\cdot10^{-6},
 \qquad
 \mathcal R_{\mathrm{tr},\infty}=-1.84\cdot10^{-5},
\end{equation}
то есть перенос восстановлен как нулевая мода в пределах дискретизационной
погрешности.

\section{Судьба отрицательных кандидатов}

После единственной знаковой поправки прежние критические блоки становятся
положительными. Их экстраполированные нижние уровни равны
\begin{equation}
 \lambda_{(1,-1),\infty}=3.68989479,
 \qquad
 \lambda_{(1,-2),\infty}=3.72136029.
\end{equation}
В окне трёх характеров $c=-1,0,1$ и меток $-3\le j\le3$ отрицательного
низколежащего уровня не найдено. Сопряжённые сектора совпадают с остатком
$2.71\cdot10^{-14}$. Следовательно, значения $-0.02386058$ и
$-0.01112057$ были артефактом несогласованного знака фоновой калибровки, а
не физическими направлениями распада нити.

\begin{nogocriterion}
Исправление отрицательных кандидатов не является сертификатом полной
устойчивости: вычислено конечное окно $-3\le j\le3$ и низколежащая часть
спектра. Пока не закрыт полный высокоугловой хвост трёх скрученных секторов,
устойчивость прямой нити и рождение материи остаются открытыми.
\end{nogocriterion}

\section{Следующий различающий тест}

Нулевая гипотеза: центробежная часть полного $Q+T+B$ оператора обеспечивает
положительную коэрцитивность при достаточно больших $|j|$, а конечный мост
от $|j|=4$ до аналитического порога не содержит отрицательных уровней.

Стоп-критерий: любой отрицательный блок после точной переносной калибровки
закрывает устойчивость прямой нити. Если конечный мост и аналитический хвост
положительны, линейная устойчивость полного прямого вихря будет закрыта.

\begin{protocol}
\begin{enumerate}
 \item полный стационарный фон $(K,a,b,q)$: \textbf{использован};
 \item локализованная ошибка: \textbf{знак материального члена фоновой калибровки};
 \item число изменённых членов оператора: \textbf{один};
 \item переносная нулевая мода: \textbf{восстановлена};
 \item перекрытие на сетке $280$: \textbf{$0.99999986$};
 \item прежние отрицательные пары: \textbf{сняты как артефакт};
 \item проверенное окно: \textbf{$c=-1,0,1$, $-3\le j\le3$};
 \item полный высокоугловой хвост: \textbf{открыт};
 \item рождение материи: \textbf{не закрыто};
 \item следующий гейт: \textbf{коэрцитивность полного высокоуглового хвоста}.
\end{enumerate}
\end{protocol}

Машинный сертификат сохранён в
\path{s2t_v6_bosonic_defect_full_tensor_translation_calibration_gate_results.json}.

\begin{thebibliography}{9}
\bibitem{translation-calibration-gustafson}
S.~Gustafson, I.~M.~Sigal,
\emph{The Stability of Magnetic Vortices}, arXiv:math/9904158.

\bibitem{translation-calibration-alonso}
J.~M.~Alonso-Izquierdo, W.~Garc\'ia Fuertes, J.~Mateos Guilarte,
\emph{Dissecting zero modes and bound states on BPS vortices in
Ginzburg--Landau superconductors}, arXiv:1602.09084.
\end{thebibliography}